\documentclass[11pt,letterpaper]{article}
\usepackage{fontspec}
\usepackage{tocloft}
\usepackage{multicol}
\usepackage{nicefrac}
\usepackage{url}
\usepackage[left=1.3in,right=1.3in,top=1in,bottom=1in]{geometry}
\setmainfont[Scale=1.0,AutoFakeBold=1.5,AutoFakeSlant=0.3]{Doves Type}

% Prevent page breaks within list items
\widowpenalty=10000
\clubpenalty=10000
\interlinepenalty=500

\title{Chicken Consomm\'e (Freeze-Thaw)}
\author{}
\date{}

\begin{document}

\maketitle
\thispagestyle{empty}

\begin{quote}\textit{Roast chicken stock is frozen into ice cubes and left to thaw slowly in a strainer in the refrigerator, where the stretched gelatin mesh lets the clear essence drip through and holds back the fat and the hazy proteins. What drips is colourless, twice as concentrated, and will not set. Makes about 2\nicefrac{1}{2}~cups.}\end{quote}

\section*{Ingredients}
\setlength{\columnsep}{20pt}
\noindent
    Chicken stock (60-Minute Roast Chicken Stock) \dotfill 1~batch (5~cups) \\
    Kosher salt \dotfill to taste (about \nicefrac{1}{8}~tsp. per cup) \\
    Unflavored powdered gelatin, optional, to set it \dotfill 1\nicefrac{1}{2}~tsp. \\

\section*{Directions}

\noindent
Chill the \textbf{stock} and skim any cold fat from its surface ---
Fill the ice cube trays ---
Set a fine-mesh strainer over a bowl, and clear room for it in the refrigerator

\begin{enumerate}
    \item Pour the \textbf{chicken stock} into the trays and freeze until the cubes are solid, at least \textit{8~hours} or overnight. Small cubes freeze fast, and fast freezing is what builds the ice crystals that stretch the gelatin mesh.

    \item Tip the frozen cubes into a fine-mesh strainer set over a bowl and put the whole thing in the refrigerator. Then walk away: the cubes want \textit{12--24~hours} to thaw slowly, and the cold is doing the work. Fat stays solid at refrigerator temperature and cannot drip through, and the hazy proteins are too large to fit the stretched mesh; warm the stock and the fat melts and the mesh relaxes, and the haze goes straight through.

    \item Let it drip until the cubes have collapsed and the drips stop. What has collected below should be colourless and clear, roughly half the volume you started with and twice as concentrated. Hold it up to the light; it should look like pale whiskey, not like stock.

    \item Leave the residue alone rather than pressing it, and keep it: it is intensely gelatinous, so a spoonful melted into a pan sauce gives the sauce body.

    \item Taste the \textbf{consomm\'e} and season it with \textbf{kosher salt}, a pinch at a time, about \nicefrac{1}{8}~tsp. per cup; it needs far less than you expect, since concentrating the stock doubled everything already in it. Bottle it, or serve it in small cups.

    \item To serve it as a jelly, bloom the \textbf{powdered gelatin} in 2~Tbsp. cold water for \textit{5~minutes}, then warm 1~cup of the \textbf{consomm\'e} with the bloomed gelatin until it dissolves, stir it back through the rest, and chill for \textit{4~hours}. Without it the consomm\'e stays liquid, because the gelatin that would have set it never made it through the strainer.
\end{enumerate}

\newpage

{\scriptsize
\setlength{\columnsep}{20pt}
\setlength{\multicolsep}{6pt}
\begin{multicols}{2}
\setlength{\parindent}{0pt}
\setlength{\parskip}{2pt}
\setlength{\itemsep}{0pt}
\setlength{\parsep}{0pt}

\subsection*{Yield}
\begin{itemize}
    \item Makes about 2\nicefrac{1}{2}~cups, plus roughly an equal volume of gelatinous residue
\end{itemize}

\subsection*{Equipment Required}
\begin{itemize}
    \item Three standard ice cube trays, or two trays plus a small freezer container (about 5~cups in all)
    \item Fine-mesh strainer, and a bowl or 4-cup measuring cup to catch the drips
    \item Refrigerator space for an overnight thaw
    \item Small saucepan and measuring spoons (only if setting it with \textbf{gelatin})
    \item Funnel, and a bottle or small jars for storage
\end{itemize}

\subsection*{Mise en Place}
\begin{itemize}
    \item Chill the \textbf{stock} completely and lift any solid fat from its surface
    \item Fill the trays with the full batch and make room in the freezer
    \item Set the strainer over the bowl and clear a shelf in the refrigerator where it can sit undisturbed
    \item If you are setting it: bloom \textit{1\nicefrac{1}{2}~tsp.} \textbf{powdered gelatin} in 2~Tbsp. cold water
\end{itemize}

\subsection*{Ingredient Tips}
\begin{itemize}
    \item \textbf{Stock}: it has to be gelatin-rich, which is why this follows the pressure-cooked stock. From a stock cooked too long at pressure, or a thin simmered one, the consomm\'e drips cloudy and tastes weak
    \item \textbf{Powdered gelatin}: 0.5--1\% by weight is the range for a soft jelly, and one 7~g packet sets about 3~cups of liquid. Sheet gelatin works too, about 4 sheets to a packet
    \item \textbf{Kosher salt}: the consomm\'e is twice as concentrated as the stock, so it will not take the same measure---taste before you salt
\end{itemize}

\subsection*{Preparation Tips}
\begin{itemize}
    \item Cover the bowl loosely; a consomm\'e sitting under a strainer for a day will take on whatever else is in the refrigerator
    \item Do not stir, press, or prod the slush to hurry it. Disturbing what is draining carries the haze down along with the clear liquid
    \item Chip any fat cap off the cubes before they thaw. The mesh would catch the fat anyway; this simply keeps the consomm\'e fat-free
    \item A stalled drip is normal---the gelatin holds on to some of the water by design, and roughly half the volume is the yield
\end{itemize}

\subsection*{Make Ahead \& Storage}
\begin{itemize}
    \item Keeps refrigerated \textit{4--5~days}, frozen up to \textit{3~months}, and is best within \textit{2--3~days}: the roast aroma fades long before it spoils
    \item It will not set on its own, and that is expected, not a sign the method failed
    \item The residue keeps like the stock---\textit{1~week} refrigerated, \textit{3~months} frozen. Freeze it in cubes, one per pan sauce
    \item Freeze the stock rather than the consomm\'e if you need to keep it longer than a week; the filtration is quick to repeat
\end{itemize}

\subsection*{Serving Suggestions}
\begin{itemize}
    \item Serve it hot in small cups at the start of a meal, with nothing in it to get in the way
    \item Finely diced carrot, leek, and celery added just before serving make it consomm\'e brunoise
    \item Use it where intensity matters: to finish a sauce, enrich a risotto, or loosen a pur\'ee. The residue gives a deglazed pan the body of a demi-glace
\end{itemize}

\subsection*{Source}
\begin{itemize}
    \item Adapted from Chris Young (Combustion Inc.; formerly The Fat Duck), via YouTube \textit{\$6 Michelin Stock in 60 Minutes (Costco Hack)}: \url{https://www.youtube.com/watch?v=3k20zFlbFfE}
\end{itemize}

\end{multicols}
}

\end{document}
